% Authority: exact-scope leaf 21; full PDF page 100; printed p.83. Translation of direct transcription.
% U:p083.u01 | continuation
\noindent
If one compares what has just been proved with the result $(\theta)$, obtained above,
one will see that $a$ is or is not a biquadratic residue with respect to $p$, according
as one has:
\[
\leg{\psi+t}{a}=\leg{\psi-t}{a}=\leg{t}{a},
\]
or:
\[
\leg{\psi+t}{a}=\leg{\psi-t}{a}=-\leg{t}{a}.
\]
% U:p083.u02 | continuation
The equation $p=t^2-au^2$ shows that, if one denotes by $\chi$ either
of the two roots of the congruence $\chi^2\equiv p\md{a}$, one has $t\equiv\chi$ or $t\equiv-\chi$
$\md{a}$. It follows from this, $a$ being of the form $4n+1$, that $\leg{t}{a}=\leg{\chi}{a}$.
On the other hand, since by virtue of what precedes $\leg{\psi+t}{a}=\leg{\psi-t}{a}$, one sees that one
always has:
\[
\leg{\psi+\chi}{a}=\leg{\psi+t}{a}=\leg{\psi-t}{a}.
\]
% U:p083.u03 | continuation
One concludes from this, having regard to what was said only a moment ago,
that the number $a$ will be a biquadratic residue with respect to $p$ when one has
$\leg{\chi+\psi}{a}=\leg{\chi}{a}$ or, which is the same thing, when
$\leg{\chi(\chi+\psi)}{a}=1$, and that
$a$ will be a biquadratic non-residue with respect to $p$ when one has $\leg{\chi+\psi}{a}=-\leg{\chi}{a}$,
or, which comes to the same thing, when $\leg{\chi(\chi+\psi)}{a}=-1$.
% U:p083.u04 | paragraph-start
\par\indent
This result concerns the case in which $\varphi$ is not divisible by $a$. It remains to
examine the case of $\varphi$ divisible by $a$. Treating this case in a similar
way, one finds that, when it holds, one of the numbers $t+\psi$, $t-\psi$ (we shall
denote it by $t\pm\psi$) is divisible by $a$ and that, as in the first case,
the other, $t\mp\psi$, is such that:
\[
\leg{t\mp\psi}{a}=(-1)^{\frac{p-1}{4}+\frac{t-1}{2}},
\]
or, which comes to the same thing, on adding to $t\mp\psi$ the number $t\pm\psi$, a multiple of $a$:
\[
\leg{2t}{a}=(-1)^{\frac{p-1}{4}+\frac{t-1}{2}}.
\]
% U:p083.u05 | paragraph-start; continues on 84
\par\indent
Comparing this with what was proved above $(\theta)$, one will find
that, in the case of $\varphi$ divisible by $a$, $a$ is or is not a biquadratic residue
with respect to $p$, according as one has $\leg{2}{a}=1$, or $\leg{2}{a}=-1$, or, which is the
